\documentclass[10pt,a4paper]{amsart}
\usepackage[margin=19mm]{geometry}
\usepackage{amsmath,amssymb,mathtools}
\usepackage[scr=rsfs,cal=euler]{mathalfa}
\usepackage{xfrac}
\usepackage[unicode,hidelinks,bookmarks=false]{hyperref}
\newcommand{\ie}{i.e.\ }
\newcommand{\cf}{cf.\ }
\newcommand{\resp}{respectively\ }
\newcommand{\Subsection}{\subsection}
\providecommand{\nobreakdash}{\nobreak\hbox{-}}
\setlength{\emergencystretch}{2em}
\multlinegap0pt
\usepackage{bm}
\usepackage{bbm}
\usepackage{dsfont}
\usepackage{xspace}
\usepackage{cancel}
\usepackage{mathtools}
\usepackage{amsthm}
\makeatletter
\@ifundefined{proposition}{\newtheorem{proposition}{Proposition}}{}
\makeatother
\newcommand{\T}[1][]{\mathbb{T}^#1}
\newcommand{\Z}[1][]{\mathbb{Z}^#1}
\newcommand{\dK}[1]{\delta_{K}^{#1}}
\newcommand{\laplac}[1]{\Delta_{#1}}
\newcommand{\la}{\lambda}
\newcommand{\supp}[1]{\text{supp}\, #1}
\title[Lossless Strichartz estimates on rectangular tori over short]{Lossless Strichartz estimates on rectangular tori over short time intervals}
\author{Connor Quinn}
\date{Source: arXiv:2601.09895v2 (CC BY 4.0)}
\begin{document}
\maketitle

\noindent\emph{Source and licence.} Lossless Strichartz estimates on rectangular tori over short time intervals. Version arXiv:2601.09895v2 (CC BY 4.0), distributed under the Creative Commons Attribution 4.0 International licence, as recorded in the arXiv licence metadata for this exact version and shown on the abstract page https://arxiv.org/abs/2601.09895v2. Full rights notice: licenses/arxiv-2601.09895-CC-BY-4.0.txt.\par\smallskip

\noindent\textit{Editorial scope.} Counted unit: the Proposition whose source label is \texttt{None} in the article (source0.tex lines 792--801, source label \texttt{None}), delivered together with the article's own complete original proof of it (source0.tex lines 803--865, one \texttt{proof} environment of 63 lines). No mathematics is written, repaired or re-derived: statement and proof are the article's own text line for line. Required context retained uncounted: stationary phase estimate for standard kernel (lines 763--765). Every definition, display and label of those lines is retained unchanged. Omitted from this package and not redistributed: 1--762, 766--791, 802--802, 866--1351. Nothing inside a retained excerpt is omitted or altered: every retained body is delivered exactly as the article's lines are written, so no omission receipt of any kind accompanies this package. Citations: the retained excerpts carry no citation command at all, so no bibliography entry of the article is transcribed and no reference is redistributed. Reference adaptation: none was needed and none was made; every reference inside the delivered unit resolves inside it, so no unresolved reference or citation is produced. Printed slips kept literal: no printed word, symbol, hypothesis or number of the counted statement or of its proof was altered. Layout and class: the article's own class is \texttt{amsart} with packages including amsmath, amsfonts, amsthm, euscript, amssymb, graphicx, mathrsfs, enumitem, indentfirst, csquotes, bbm; the wrapper is the lane's pinned \texttt{amsart} editorial wrapper at 10pt on A4, which restates the theorem-like environments the retained text needs and adds the article's own macro definitions that the retained text needs (\texttt{\textbackslash T}, \texttt{\textbackslash Z}, \texttt{\textbackslash dK}, \texttt{\textbackslash la}, \texttt{\textbackslash laplac}, \texttt{\textbackslash supp}), with extra wrapper packages (bm, bbm, dsfont, xspace, cancel, mathtools). The delivered numbering is the unit's own. Assets: no retained span contains a figure environment, a graphics-inclusion command, a PDF-page inclusion, a table or a TikZ picture; no image file is copied into this package, the package declares no asset dependency and no image file is redistributed with it. Licence: version arXiv:2601.09895v2 of this article is distributed under the Creative Commons Attribution 4.0 International licence, as recorded in the arXiv licence metadata for this exact version and shown on the abstract page https://arxiv.org/abs/2601.09895v2. A full rights notice is delivered at licenses/arxiv-2601.09895-CC-BY-4.0.txt. Characters: every retained excerpt is pure ASCII, so the delivered engine, XeLaTeX with its default Latin Modern text font, reports no missing character.
\begin{equation} \label{stationary phase estimate for standard kernel}
    \int \phi^2(\xi/\la)e(\la^{-1}(t-s)|\xi|^2 +  (x-y) \xi)d \xi = O\big(\la^{\frac{1}{2}}|t-s|^{-\frac{1}{2}}\big).
\end{equation}

\begin{proposition}
Let $\kappa > 0$ and suppose $\delta \geq \la^{-1 + \kappa}$. Then for $x,y \in \T{}$ we have
\begin{equation} 
     \big|\eta(\delta t)\eta(\delta s)\big( \chi_{I}^{2}(D_x)\phi^2(D_x / \la)e^{-i\la^{-1}(t-s)\laplac{\T{}}} \big)(x,y)\big| \lesssim \la^{\frac{1}{2}}|t-s|^{-\frac{1}{2}}\big(1 + \dK{m}|t-s|\big).
\end{equation} 
Thus, when $m = K$ (so that $\dK{m} = \delta$), if $|t-s| \leq 2\delta^{-1}$ we have 
\begin{equation}
    \big|\eta(\delta t)\eta(\delta s)\big( \chi_{I}^{2}(D)\phi^2(D / \la)e^{-i\la^{-1}(t-s)\laplac{\T{}}} \big)(x,y) \big| \lesssim \la^{\frac{1}{2}}|t-s|^{-\frac{1}{2}}.
\end{equation}
\end{proposition}

\begin{proof}
First, fix $x,y \in \T{}$ and note that $\partial_{\xi}^{k}\chi_{I}(\xi) = O\big( (\la \dK{m})^{-k}\big)$. Also, this function is supported in an interval of length $\sim \la \dK{m}$ so  $\check{\chi}_{I}(x)$ =  $O\big( (\la \dK{m})(1+\la \dK{m}|x|)^{-N}\big)$ for any $N \geq 0$. It then follows from the Poisson Summation formula that the kernel of our operator satisfies
$\chi_I(D_x)(x,y) = O\Big(\la \dK{m}\big(1 + \la \dK{m}|x- y|\big)^{-N}\Big)$ for any $N \geq 0$. It follows that
\begin{equation} \label{kernel bounds for angular cutoff multipliers}
    \int_{\mathbb{R}}|\check{\chi}_{I}(z)|dz, \ \sup_{y \in \mathbb{T}}\int_{\mathbb{T}} |\chi_{I}(D_x)(x,y)|dx , \ \sup_{x \in \mathbb{T}}\int_{\mathbb{T}} |\chi_{I}(D_x)(x,y)|dy \leq C.
\end{equation}
Thus, it follows from Proposition 6.1 that we obtain the desired bounds in Proposition 6.2 if $|t-s| \leq 1$.
\par
Also, (\ref{kernel bounds for angular cutoff multipliers}) implies that if we let 
\begin{equation}
    K^{\la}_{I}(x,t;y,s) = \eta(\delta t)\eta(\delta s)\big(\chi_{I}(D_x)\phi^{2}(D_x / \la)e^{-i\la^{-1}(t-s)\laplac{\T{}}}\big)(x,y)
\end{equation}
then it suffices to see that this kernel satisfies the desired bounds if $|t-s| \geq 1$. Also, recall that
\begin{equation} 
    \chi_I(\xi) = \chi\big((\la \dK{m})^{-1}(\xi - c)\big)
\end{equation}
and
\begin{equation} \label{support property of chi_I}
    \supp{\chi_I} \subset 2I = [c -2\la\dK{m}, c+2\la \dK{m} ]
\end{equation}
\par
If we let $K^{\la}(x,t;y,s)$ be the kernel of the multiplier operator $\eta(\delta t)\eta(\delta s)\phi^2(D_x/\la)e^{-i\la^{-1}(t-s)\laplac{\T{}}}$ then using the Poisson Summation formula and writing $\chi_{I}$ in terms of its inverse Fourier transform we have
\begin{equation} \label{K^(la, theta) Poisson summation 1}
\begin{split}
    K^{\la}_{I} (x,t;y,s) &= \eta(\delta t) \eta(\delta s)\frac{1}{2 \pi} \sum_{j \in \Z{}}\int_{\mathbb{R}} \chi_{I}(\xi)\phi^2(\xi/\la)e((x-y-j)\xi + \la^{-1}(t-s)|\xi|^2)d\xi\\
    &= \eta(\delta t) \eta(\delta s) \frac{1}{2\pi}\int_{\mathbb{R}} \check{\chi}_{I} (x - z)K^{\la}(z,t; y, s) dz.
\end{split}
\end{equation}
Further, if we repeat the arguments from the proof of Proposition 6.1 then we see that for $|t-s| \in [2^{j-1}, 2^j]$ we have
\begin{equation} \label{K^la poisson summation}
    K^{\la}(z,t; y, s) = \sum_{j \in \Z{}}\tilde{K}^{\la}(z,t;y+j,s) + O(\la^{-N})
\end{equation}
where
\begin{equation} \label{tilde K^la definition} 
    \tilde{K}^{\la}(z,t; y, t) = \frac{1}{2\pi} \int_{\mathbb{R}} \phi^2(\xi/\la)e(\la^{-1}(t-s)|\xi|^2 + (z-y) \xi)d\xi.
\end{equation}
As before, the sum in (\ref{K^la poisson summation}) is over the $j \in \Z{}$ such that $|z - (y + j)| \lesssim |t-s|$ due to Huygen's principle (in particular, the sum is finite). It follows that 
\begin{equation} \label{K^(la, theta) Poisson summation 2}
    K^{\la}_{I} (x,t;y,s) = \eta(\delta t) \eta(\delta s) \sum_{j \in \Z{}}\int_{\mathbb{R}} \check{\chi_{I}} \ (x - z)\tilde{K}^{\la}(z,t; y+ j, s) dz + O\big(\la^{-N}\big)
\end{equation}
We recall the estimate (\ref{stationary phase estimate for standard kernel}), which says that 
\begin{equation} \label{stationary phase bound for K^la}
    \big| \tilde{K}^{\la}(z,t;y,s)\big| \lesssim \la^{\frac{1}{2}}|t-s|^{-\frac{1}{2}}.
\end{equation}
It follows from (\ref{kernel bounds for angular cutoff multipliers}) and (\ref{stationary phase bound for K^la}) that for each fixed $j$ we have
\begin{equation}
    \Big|\int_{\mathbb{R}} \check{\chi}_{I} \ (x - z)\tilde{K}^{\la}(z,t; y+ j, s) dz\Big| \lesssim \la^{\frac{1}{2}}|t-s|^{-\frac{1}{2}}
\end{equation}
Thus, in order to finish the proof of Proposition 6.2 it suffices to see that that number of non-negligible summands in (\ref{K^(la, theta) Poisson summation 2}) is $O\big(\dK{m}|t-s|\big)$.
\par
Note that 
\begin{equation}
    \int_{\mathbb{R}} \check{\chi}_{I} \ (x - z)\tilde{K}^{\la}(z,t; y+ j, s) dz = \int_{\mathbb{R}} \chi_{I}(\xi)\phi^2(\xi/\la)e(( x-y-j) \xi + \la^{-1}(t-s)|\xi|^2)d\xi.
\end{equation}
Recalling (\ref{support property of chi_I}) we can integrate by parts to see that this integral is $O(\la^{-N})$ for all $N \geq 0$ if $j$ is not contained in the set
\begin{equation}
    \{ j \in \Z{} \ : \ |x- y - j|\in \big[2|t-s|(\la^{-1}c - C_0\dK{m}), 2|t-s|(\la^{-1}c + C_0\dK{m})\big] \}
\end{equation}  
for fixed $x,y,t$ and $s$, where $C_0$ is a large fixed constant. This is an interval of length $\sim \dK{m}|t-s|$ so we have that
\begin{equation}
    \#\{ j \in \Z{} \ : \ |x- y - j|\in \big[2|t-s|(\la^{-1}c - C\dK{m}), 2|t-s|(\la^{-1}c + C\dK{m})\big] \} = O\big(\dK{m}|t-s|\big).
\end{equation} Thus, the proof of Proposition 6.2 is complete.
\end{proof}

\end{document}
